% Part: first-order-logic
% Chapter: axiomatic-proofs
% Section: identity

\documentclass[../../../include/open-logic-section]{subfiles}

\begin{document}

\olfileid{fol}{axd}{ide}

\olsection{له \usetoken{S}{identity} سره \usetoken{P}{derivation}}

په !!{derivation}s کښې د~$\eq$ د شاملولو دپاره يوازې نوې بديهي
شېماوې ورزياتوو۔ د~!!{derivation} او~$\Proves$ تعريف هماغسې پاتې
کېږي؛ يوازې نوي بديهي اصول هم منو۔

\begin{defn}[د !!{identity} بديهي اصول]
\begin{align}
\ollabel{ax:id1} & \eq[t][t], \\
\ollabel{ax:id2} & \eq[t_1][t_2] \lif (!B(t_1) \lif
  !B(t_2)),
\end{align}
د هر تړلي ترم~$t$، $t_1$ او~$t_2$ دپاره۔
\end{defn}

\begin{prop}\ollabel{prop:sound}
  بديهي اصول~\olref{ax:id1} او~\olref{ax:id2} معتبر دي۔
\end{prop}

\begin{proof}
  تمرين۔
\end{proof}

\begin{prob}
\olref[fol][axd][ide]{prop:sound} ثابت کړئ۔
\end{prob}

\begin{prop}\ollabel{prop:iden1}
 د هر تړلي ترم~$t$ او سټ~$\Gamma$ دپاره
 $\Gamma \Proves \eq[t][t]$۔
 \textbf{د ژباړې د نسخې سمون (OLAX-011):} سرچينې دلته «هر ترم»
 ويلي وو، خو د يووالي لومړے بديهي اصل يوازې پر تړلو ترمونو شېما
 ده؛ هماغه
 لازم قيد بېرته څرګند شو۔
\end{prop}

\begin{prop}\ollabel{prop:iden2}
  که $\Gamma \Proves !A(t_1)$ او $\Gamma \Proves
  \eq[t_1][t_2]$ وي، نو $\Gamma \Proves !A(t_2)$، په دې شرط چې
  دواړه ښودل شوي ترمونه تړلي وي۔
  \textbf{د ژباړې د نسخې سمون (OLAX-012):} سرچينې په دې قضيه کښې
  د تړلو ترمونو شرط نۀ و بيان کړے، سره له دې چې ثبوت ئې نېغ په
  نېغه د يووالي دويم بديهي اصل کاروي او هغه شېما يوازې تړلي
  ترمونه مني۔
\end{prop}

\begin{proof}
دا !!{formula}
\[
(\eq[t_1][t_2] \lif (!A(t_1) \lif !A(t_2)))
\]
د~\olref{ax:id2} يوه بېلګه ده۔ پايله د~\MP په دوه پلي کولو سره
ترلاسه کېږي۔
\end{proof}

\end{document}
